\chapter{受身形 — Voz Pasiva}
\unidadheader{28}{受身形}{Voz Pasiva}{Expresar acciones recibidas}

\begin{tcolorbox}[vocabbox, title={📌 Vocabulario Nuevo}]
\begin{tabularx}{\linewidth}{llX}
\toprule
\textbf{Japonés} & \textbf{Español} & \textbf{Tipo} \\
\midrule
\vocabitem{怒る}{おこる・おこられる}{enojarse / ser regañado}{v. gr.1}
\vocabitem{ほめる}{—}{elogiar}{v. gr.2}
\vocabitem{叱る}{しかる}{regañar / reprender}{v. gr.1}
\vocabitem{助ける}{たすける}{ayudar / rescatar}{v. gr.2}
\vocabitem{守る}{まもる}{proteger}{v. gr.1}
\vocabitem{批判する}{ひはんする}{criticar}{v. irr.}
\vocabitem{評価する}{ひょうかする}{evaluar / valorar}{v. irr.}
\vocabitem{尊敬する}{そんけいする}{respetar}{v. irr.}
\vocabitem{信頼する}{しんらいする}{confiar en}{v. irr.}
\vocabitem{傷つける}{きずつける}{herir / lastimar}{v. gr.2}
\bottomrule
\end{tabularx}
\end{tcolorbox}

\begin{tcolorbox}[phrasebox, title={⚙️ Conjugación: Forma pasiva}]
\begin{tabular}{llll}
\toprule
\textbf{Grupo} & \textbf{Regla} & \textbf{Original} & \textbf{Pasiva} \\
\midrule
1 (う→あ段+れる) & う→あれる & \ruby{読}{よ}む & \ruby{読}{よ}まれる \\
2 (〜る → 〜られる) & る→られる & \ruby{食}{た}べる & \ruby{食}{た}べられる \\
3 & memorizar & する / くる & される / こられる \\
\bottomrule
\end{tabular}
\end{tcolorbox}

\subsection*{🗣️ Frases Clave}

\frase{\ruby{先生}{せんせい}にほめられました。}{El profesor me elogió. (me fue elogiado)}
\frase{あの\ruby{映画}{えいが}は\ruby{世界中}{せかいじゅう}で\ruby{見}{み}られています。}
  {Esa película es vista en todo el mundo.}
\frase{\ruby{財布}{さいふ}を\ruby{盗}{ぬす}まれました。}{Me robaron la cartera.}

\subsection*{⚙️ Gramática}

\patron{Pasiva directa}
  {A が B に V-られる (A recibe la acción de B)}
  {\ej{\ruby{私}{わたし}は\ruby{先生}{せんせい}にほめられました。}
    {El profesor me elogió (lit. fui elogiado por el profesor).}
  \ej{この\ruby{本}{ほん}は\ruby{多く}{おおく}の\ruby{人}{ひと}に\ruby{読}{よ}まれています。}
    {Este libro es leído por muchas personas.}}

\patron{Pasiva de molestia (indirect passive)}
  {A が B に C を V-られる (molesta al sujeto)}
  {\ej{\ruby{電車}{でんしゃ}の\ruby{中}{なか}で\ruby{足}{あし}を\ruby{踏}{ふ}まれました。}
    {Me pisaron el pie en el tren. (me fue pisado, molestia)}
  \ej{\ruby{雨}{あめ}に\ruby{降}{ふ}られて、\ruby{濡}{ぬ}れました。}
    {Me cayó la lluvia y me mojé. (molestia de la lluvia)}}

\begin{tcolorbox}[ejerciciobox, title={📝 Ejercicios Rápidos}]
\begin{enumerate}
  \item Convierte a pasiva: 先生が学生をほめる → ?
  \item ¿Cuándo se usa la pasiva de molestia?
  \item Traduce: \ruby{子}{こ}どもの\ruby{時}{とき}、よく\ruby{母}{はは}に\ruby{叱}{しか}られました。
\end{enumerate}
\vspace{0.4em}
\textbf{Respuestas:}
1) \ruby{学生}{がくせい}は\ruby{先生}{せんせい}にほめられた。\quad
2) Cuando la acción de otro te afecta negativamente.\quad
3) Cuando era niño/a, mi mamá me regañaba (era regañado) seguido.
\end{tcolorbox}

\begin{tcolorbox}[kanjibox, title={🀄 Kanjis de la Unidad}]
\begin{tabularx}{\linewidth}{cllXX}
\toprule
\textbf{Kanji} & \textbf{On} & \textbf{Kun} & \textbf{Significado} & \textbf{Vocab} \\
\midrule
\kanjirow{怒}{ど}{おこ・いか}{enojarse}{\ruby{怒}{おこ}る / \ruby{怒り}{いかり}}
\kanjirow{助}{じょ}{たす}{ayudar / rescatar}{\ruby{助}{たす}ける / \ruby{助言}{じょげん}}
\kanjirow{守}{しゅ}{まも・もり}{proteger / guardar}{\ruby{守}{まも}る / \ruby{守備}{しゅび}}
\kanjirow{批}{ひ}{—}{criticar}{\ruby{批判}{ひはん} / \ruby{批評}{ひひょう}}
\kanjirow{評}{ひょう}{—}{evaluar}{\ruby{評価}{ひょうか} / \ruby{評判}{ひょうばん}}
\bottomrule
\end{tabularx}
\end{tcolorbox}
